\chapter{ОПИСАНИЕ ПРОГРАММНОГО ПРОТОТИПА}\label{ch:ch7}

В предыдущих главах были последовательно сформулированы доменная модель
системы автономных мобильных роботов с поддержкой периферийных вычислений,
VCG-подобный слой распределения задач, многоагентный контур выбора
операционных режимов периферийных узлов, гибридный метод AMPLE-AMR, а также
доменная параметризация роботизированного склада для имитационного
моделирования. Для
перехода от формальной постановки к экспериментальной проверке был создан
программный прототип, реализующий воспроизводимую имитационную среду для роботизированного склада.

Цель настоящей главы состоит в описании архитектуры этого прототипа,
принятых программных решений и механизма воспроизводимого получения
результатов. Описание ориентировано не на пользовательский интерфейс, а на
исследовательский контур: как в коде представлены роботы, вычислительные
задачи, периферийные узлы, сетевые связи, слои распределения и контур
обучения, а также как фиксируются параметры сценариев и экспортируются
артефакты для последующего анализа.

\section{Назначение и требования к прототипу}

Программный прототип разрабатывался как исследовательский инструмент для
решения трёх задач:
\begin{enumerate}
	\item воспроизводимое моделирование локальной AMR-системы, в которой
	роботы формируют вычислительные задачи и используют общую периферийную
	инфраструктуру;
	\item сравнение нескольких методов принятия решений в одинаковых
	сценариях, включая фиксированные базовые методы, гибридный метод
	AMPLE-AMR и кластерный вариант C-AMPLE-AMR;
	\item автоматическая генерация артефактов эксперимента в формате, пригодном
	для последующего включения в диссертационный текст: CSV-файлов,
	графиков и LaTeX-таблиц.
\end{enumerate}

При проектировании прототипа были зафиксированы следующие инженерные
требования.

\textbf{Во-первых}, все источники случайности должны контролироваться seed.
Это относится к генерации задач, начальному размещению роботов, параметрам
связи, деградации сети, отказам узлов и обучению многоагентной политики.

\textbf{Во-вторых}, параметризация сценариев не должна быть рассеяна по коду.
Для этого масштабы складов, классы задач, типы периферийных узлов,
операционные режимы и параметры сценариев вынесены в конфигурационный YAML.

\textbf{В-третьих}, обучаемый контур не должен подменять собой экономический
механизм распределения задач. Поэтому реализация QMIX отвечает только за
выбор операционных режимов узлов, а не за обучение ставок, оценок ценности
или стратегии формирования заявок.

\textbf{В-четвёртых}, весь контур эксперимента должен запускаться через единую
точку входа и поддерживать режимы обучения, оценки, короткого проверочного
запуска, генерации графиков и экспорта LaTeX-таблиц.

\section{Общая архитектура программного прототипа}

Прототип реализован как Python-пакет, состоящий из нескольких логических
подсистем:
\begin{enumerate}
	\item конфигурационный слой;
	\item доменная модель роботов, задач, точек доступа и периферийных узлов;
	\item имитационная среда и сетевая модель;
	\item слои распределения задач;
	\item контур многоагентного обучения и исполнения;
	\item подсистема управления экспериментами и экспорта результатов;
	\item подсистема тестирования.
\end{enumerate}

Такое разбиение отражает архитектурную декомпозицию метода AMPLE-AMR.
Доменная и сетевая модели отвечают за формирование реалистичного состояния
складской системы. Слои распределения реализуют экономический механизм
назначения задач. Многоагентный контур отвечает за адаптивный выбор режимов
узлов. Модуль управления экспериментами связывает эти блоки в единый процесс,
обеспечивая воспроизводимый запуск и сохранение артефактов.

\section{Структура пакета и распределение ответственности}

Реализация организована как пакет \texttt{ample\_amr}. Основные модули
распределены следующим образом.

\textbf{Модуль конфигурации} загружает сценарии, размеры складов, параметры
классов задач и параметры обучения. Он также преобразует YAML-описание в
типизированные объекты конфигурации, используемые во всех остальных частях
системы.

\textbf{Модуль доменной модели} содержит классы \texttt{Task},
\texttt{Robot}, \texttt{EdgeNode}, \texttt{AccessPoint} и вспомогательные
структуры для описания резервирования и локального состояния узла.

\textbf{Модуль \texttt{workload}} реализует генерацию задач по индивидуальным seed роботов и
классовым распределениям. В нём формируется смесь задач обнаружения
препятствий, восприятия, локализации, перепланирования, инвентаризации и
диагностики.

\textbf{Модуль \texttt{network}} отвечает за построение графа беспроводной
достижимости между роботами, точками доступа и периферийными узлами,
оценивание задержек и пропускной способности, а также за вычисление
кластеризационной аффинности между роботами.

\textbf{Модуль \texttt{utility}} реализует функции оценки сетевого времени, времени
ожидания, времени обработки, полезности, стоимости обслуживания и
общественного благосостояния.

\textbf{Модуль \texttt{allocation}} содержит общий интерфейс распределителя задач и три
реализации: эвристическую, VCG-подобную и кластерную VCG-подобную.

\textbf{Модуль \texttt{qmix}} содержит реализацию контроллера операционных режимов
узлов: фиксированный контроллер, случайный базовый метод и QMIX-контур с буфером опыта, целевыми сетями и факторизацией функции ценности.

\textbf{Модуль \texttt{environment}} описывает имитационную среду как связку
динамики роботов, генерации задач, учёта отказов, применения режимов,
распределения задач и агрегации метрик.

\textbf{Модули \texttt{runner} и \texttt{outputs}} образуют верхний уровень управления
экспериментами. Они запускают обучение и оценку, агрегируют шаговые и
эпизодные показатели и формируют CSV-файлы, графики и LaTeX-таблицы.

\section{Конфигурационный слой и воспроизводимость сценариев}

Все основные параметры прототипа вынесены в файл конфигурации. В нём
фиксируются:
\begin{enumerate}
	\item размеры складских конфигураций Warehouse-S, Warehouse-M,
	Warehouse-M+, Warehouse-L и Warehouse-XL;
	\item характеристики классов периферийных узлов Edge-S, Edge-M и
	Edge-GPU;
	\item коэффициенты операционных режимов \texttt{safe}, \texttt{normal},
	\texttt{aggressive}, \texttt{reserve}, \texttt{critical};
	\item диапазоны параметров для каждого класса задач;
	\item параметры генерации нагрузки, деградации связи, отказов узлов и
	чувствительности к режимам;
	\item гиперпараметры обучения QMIX;
	\item значения по умолчанию для длины эпизода и длительности управляющего
	слота.
\end{enumerate}

Такое решение позволяет отделить исследовательские допущения от
реализационного кода. При добавлении нового сценария разработчик изменяет
только конфигурационный файл и при необходимости дополняет процедуру
агрегирования результатов, не переписывая базовую логику среды или
распределитель задачов.

Воспроизводимость достигается двумя средствами. Во-первых, каждый запуск
получает явный список seed, который передаётся в генераторы Python, NumPy и
PyTorch. Во-вторых, все итоговые артефакты сохраняются в файловую структуру с
фиксированными именами, что позволяет повторно использовать результаты для
построения графиков и таблиц без перезапуска симуляции.

\section{Доменная модель складской AMR-системы}

\subsection{Представление роботов}

Класс \texttt{Robot} описывает автономного мобильного робота как источник
вычислительных задач и как подвижный объект складской среды. Для каждого
робота хранятся его идентификатор, координаты, вектор скорости, максимальная
и рабочая скорости, полезная нагрузка, важность миссии и локальная сводка о
качестве сетевой связности.

В текущей версии прототипа движение робота моделируется упрощённо: на каждом
шаге робот изменяет положение внутри прямоугольной складской области и при
необходимости меняет направление. Это достаточно для анализа влияния
локальной сетевой связности и меняющихся маршрутов на распределение задач, но
не претендует на детальную физическую модель движения платформы.

\subsection{Представление вычислительных задач}

Класс \texttt{Task} является центральным объектом взаимодействия между
роботом, аукционным слоем и периферийным узлом. Задача содержит идентификатор
робота, класс задачи, момент поступления, объёмы входных и выходных данных,
нормированные требования по CPU, памяти и ускорителю, срок обработки,
приоритет, признак жёсткого дедлайна, статус исполнения, а также экономические
поля: полезность, стоимость, вклад в общественное благосостояние, платежи и
оценку внешнего эффекта.

Наличие этих полей важно по двум причинам. Во-первых, симулятор должен
поддерживать как оперативные характеристики исполнения, так и экономическую
интерпретацию распределения. Во-вторых, именно эти признаки используются для
анализа результатов и формирования локальных наблюдений узлов в
последующих шагах.

\subsection{Представление периферийных узлов и точек доступа}

Класс \texttt{EdgeNode} описывает периферийный вычислительный узел как
объект с идентификатором, типом узла, полной нормированной ёмкостью по
ресурсам, текущей загрузкой, очередью, набором исполняемых задач, текущим
операционным режимом и журналом резервирований.

Отдельный класс \texttt{AccessPoint} используется для явного представления
беспроводной инфраструктуры склада. Это решение оказалось полезным по двум
причинам:
\begin{enumerate}
	\item стало возможным моделировать трёхслойную структуру <<роботы ---
	точки доступа --- периферийные узлы>>;
	\item сетевые деградации и локальная недостижимость могут задаваться
	через региональные изменения связности, а не только через абстрактные
	штрафы в матрице стоимостей.
\end{enumerate}

\section{Имитационная среда и сетевая модель}

\subsection{Геометрия склада и сетевая связность}

Имитационная среда строит регулярную сетку точек доступа поверх складской
области. Периферийные узлы прикрепляются к отдельным точкам доступа и
получают пространственные координаты. Для каждого робота на текущем шаге
определяется ближайшая покрывающая точка доступа, после чего вычисляются
оценки задержки, пропускной способности и сетевой стоимости передачи данных
до каждого достижимого узла.

Сетевая модель различает по меньшей мере три случая:
\begin{enumerate}
	\item локальное подключение к узлу через ту же точку доступа;
	\item доступ через межточечную связность с более высокой латентностью и
	ограниченной пропускной способностью;
	\item недостижимость узла, при которой назначение запрещается.
\end{enumerate}

В сценариях деградации связи к базовой задержке добавляется дополнительный
штраф, а пропускная способность уменьшается. В результате сетевые параметры
влияют не только на стоимость обслуживания, но и на допустимость назначения
задачи в случае жёсткого дедлайна.

\subsection{Временная динамика исполнения}

На каждом управляющем шаге среда выполняет следующую последовательность:
\begin{enumerate}
	\item завершает задачи, время исполнения которых истекло;
	\item применяет события отказа и восстановления узлов;
	\item обновляет положение роботов и текущую сетевую привязку;
	\item генерирует новые задачи;
	\item применяет выбранные операционные режимы периферийных узлов;
	\item вызывает слой распределения задач;
	\item резервирует задачи на выбранных узлах;
	\item обновляет очереди, локальные наблюдения и метрики.
\end{enumerate}

Для резервирования используется календарная модель исполнения. Каждая
назначенная задача получает момент старта и завершения, а узел поддерживает
список интервалов резервирования. Это позволяет корректно учитывать
перекрывающиеся интервалы, очереди и ресурсные ограничения не только в
момент назначения, но и в течение всего горизонта ближайшего исполнения.

\section{Оценка полезности, стоимости и общественного благосостояния}

Функции \texttt{estimate\_network\_time}, \texttt{estimate\_queue\_time},
\texttt{estimate\_processing\_time}, \texttt{evaluate\_utility} и
\texttt{evaluate\_cost} реализуют численную интерпретацию формальной модели,
заданной во второй и третьей главах.

Полезность задачи зависит от:
\begin{enumerate}
	\item базовой ценности класса задачи;
	\item приоритета задачи;
	\item важности миссии робота;
	\item чувствительности к задержке;
	\item суммарной оценки задержки передачи, ожидания и обработки;
	\item штрафа за сеть и передачу данных;
	\item штрафа за нарушение дедлайна для задач с мягким сроком.
\end{enumerate}

Для задач с жёстким дедлайном среда запрещает назначение, если суммарная
оценка латентности превышает допустимый срок. Это ограничение действует до
вызова аукциона и тем самым исключает заведомо неосуществимые назначения из
множества кандидатов.

\section{Реализация слоёв распределения задач}

\subsection{Эвристический слой}

Эвристический распределитель задач поддерживает несколько политик выбора назначения.
В качестве основной базовой политики используется \texttt{greedy\_net\_welfare}:
для каждой задачи выбирается допустимое назначение с максимальным положительным
чистым вкладом в общественное благосостояние при соблюдении ограничений на
экспонируемую ёмкость узла.

\subsection{VCG-подобный слой}

Класс \texttt{VCGLikeAllocator} реализует детерминированный алгоритм
branch-and-bound. Для каждой задачи строится список допустимых кандидатов, а
затем решается задача выбора победителей под ограничениями на CPU, память и
ускоритель. Верхние оценки по остаточному положительному благосостоянию
используются для отсечения поддеревьев поиска.

Платежи роботов вычисляются по схеме Clarke pivot через повторное решение
задачи без соответствующего робота. Дополнительно рассчитывается внутренняя
компенсация узла как потеря общественного благосостояния при исключении
данного узла из пула доступных ресурсов.

\subsection{Кластерный слой}

Класс \texttt{ClusteredVCGLikeAllocator} строит граф аффинности между
роботами, где вес ребра зависит от частоты взаимодействия и сетевой стоимости.
Далее либо используется кластеризация графа, либо включается детерминированный
fallback-разбиитель. Внутри каждого кластера запускается локальный
VCG-подобный распределитель задач.

Кластеризация здесь служит именно средством
локализации распределения и анализа компромисса между накладными расходами и
потерей качества, а не самостоятельной задачей управления распределённым
протоколом.

\section{Многоагентный контур выбора режимов узлов}

В реализации обучаемого слоя каждому периферийному узлу соответствует один
агент. Пространство действий состоит из пяти режимов:
\texttt{safe}, \texttt{normal}, \texttt{aggressive}, \texttt{reserve},
\texttt{critical}. Пространство наблюдений содержит только локально доступные
или локально агрегируемые признаки: свободные доли ресурсов, длину очереди,
распределение классов задач в очереди, текущую загрузку, признак отказа,
предыдущий режим, сводку о качестве связи и итоговые показатели предыдущего
раунда распределения.

Для централизованного обучения и децентрализованного исполнения реализована
архитектура QMIX. Каждый агент использует общую сеть оценивания действий, а
глобальная командная функция ценности формируется сетью смешивания,
параметризуемой централизованным состоянием. В ходе обучения применяются:
\begin{enumerate}
	\item буфер опыта;
	\item целевые сети;
	\item epsilon-greedy exploration;
	\item периодическое обновление целевых параметров;
	\item сохранение контрольной точки для каждого масштаба склада.
\end{enumerate}

Тем самым обучаемый контур сохраняет совместимость с теоретической
постановкой главы 4 и не вмешивается в процесс формирования ценностей задач.

\section{Реализованные методы сравнения}

Прототип поддерживает несколько вариантов принятия решений, используемых в
главе 8 для проверки вклада отдельных компонентов метода.

\textbf{Fixed-Heuristic} использует фиксированные операционные режимы
периферийных узлов и эвристическое распределение задач. Этот вариант служит
простым базовым методом и позволяет оценить качество распределения без
обучаемой адаптации и без VCG-подобных платежей.

\textbf{Fixed-VCG} использует фиксированные режимы узлов и VCG-подобное
распределение задач. Этот вариант изолирует вклад аукционного слоя при
отсутствии обучаемого выбора режимов.

\textbf{QMIX-Heuristic} использует обучаемый выбор режимов периферийных узлов,
но заменяет VCG-подобный слой эвристическим распределением. Этот вариант
позволяет проверить пользу многоагентной адаптации без аукционного механизма.

\textbf{AMPLE-AMR} является полным вариантом метода: QMIX-контур выбирает
операционные режимы узлов, а VCG-подобный слой распределяет задачи в рамках
выбранных ресурсных ограничений.

\textbf{C-AMPLE-AMR} использует тот же обучаемый выбор режимов, но заменяет
глобальное распределение локальными распределениями внутри кластеров графа
инфраструктуры. В текущей версии прототипа выбор лидера кластера или
оркестратора распределения не реализуется и не используется. Кластеризация
применяется только для локализации задачи назначения и анализа накладных
расходов.

Отдельный режим \textbf{Random} может использоваться только как техническая
проверка корректности контура. Он не рассматривается как содержательный
конкурент AMPLE-AMR.

Принципиально важно, что прототип не содержит варианта <<чистый QMIX>> без слоя
распределения задач. В архитектуре работы QMIX выбирает только режимы узлов,
а непосредственное назначение задач всегда выполняет отдельный слой
распределения.

\section{Оркестрация экспериментов и экспорт артефактов}

Все сценарии запускаются через единый точку входа. Интерфейс командной строки
поддерживает:
\begin{enumerate}
	\item запуск обучающих эпизодов;
	\item оценку методов на фиксированном наборе сценариев и seed;
	\item быстрый короткого проверочного запуска для проверки контура;
	\item повторную генерацию графиков из готовых CSV;
	\item экспорт LaTeX-таблиц без повторного запуска симуляции.
\end{enumerate}

Результаты сохраняются на трёх уровнях.

\textbf{Шаговый уровень} содержит временные ряды по каждому управляющему шагу:
общественное благосостояние, задержки, доли завершённых и отклонённых задач,
загрузку узлов, накладные расходы, распределение режимов и кластерные
метрики.

\textbf{Эпизодный уровень} агрегирует показатели по эпизодам и seed.

\textbf{Итоговый уровень} усредняет результаты по seed, формирует сводные CSV,
графики и готовые LaTeX-таблицы, непосредственно используемые в главе 8.

\section{Тестирование и проверка корректности}

Для прототипа был сформирован набор модульных и интеграционных тестов.
Проверяются:
\begin{enumerate}
	\item воспроизводимость генерации задач при фиксированном seed;
	\item соблюдение диапазонов скоростей AMR;
	\item уменьшение полезности при росте задержки для
	чувствительных к задержке задач;
	\item запрет назначения задач с жёстким дедлайном при недопустимой
	оценке латентности;
	\item корректное изменение экспонируемой ёмкости при смене режима узла;
	\item соблюдение ресурсных ограничений эвристическим и VCG-подобным
	распределителями задач;
	\item согласованность знака VCG-подобных платежей роботов в тестовых сценариях с неотрицательными внешними эффектами;
	\item корректность выбора действий контроллером QMIX;
	\item отсутствие режима <<чистый QMIX>> без слоя распределения задач;
	\item запуск quick benchmark и генерация ожидаемых файлов результатов.
\end{enumerate}

Наличие такого набора тестов важно не только как инженерная практика, но и
как часть исследовательской добросовестности: ошибки в контроле seed,
ресурсных ограничений или дедлайнов непосредственно искажают результаты
экспериментальной проверки метода.


\section{Ограничения программной реализации}

Прототип является исследовательской симуляционной средой, а не промышленной
системой управления роботами. В нём сознательно упрощены компоненты, не
являющиеся предметом диссертационной работы.

Во-первых, движение роботов моделируется агрегированно и используется главным
образом для изменения сетевой связности и генерации задач. Детальное
планирование траекторий, предотвращение столкновений и низкоуровневое
управление приводами не рассматриваются.

Во-вторых, сетевая модель описывает задержки, пропускную способность и
достижимость на уровне графа. Она не воспроизводит детально физический уровень
Wi-Fi, механизмы доступа к среде и особенности конкретного сетевого
оборудования.

В-третьих, вычислительные ресурсы периферийных узлов задаются в
нормированных единицах. Это позволяет сравнивать слабые, средние и мощные
узлы, но не является прямой моделью конкретного процессора или ускорителя.

В-четвёртых, VCG-подобное распределение в прототипе реализуется как
исследовательский механизм проверки гипотезы. Для больших сценариев точное
вычисление распределения и платежей может быть заменено приближённым
решателем или кластерным вариантом.

В-пятых, качество обученной политики зависит от выбранных сценариев обучения,
диапазонов параметров и числа seed. Поэтому численные результаты главы 8
следует интерпретировать как результаты имитационной валидации в заданной
модели, а не как прямое измерение производительности промышленного склада.

\section{Выводы по главе}

В данной главе описан программный прототип, реализующий полный контур
экспериментальной проверки методов AMPLE-AMR и C-AMPLE-AMR. Показано, как
формальная постановка преобразована в воспроизводимую симуляционную среду с
явным разделением между доменной моделью, аукционным распределением, выбором
режимов узлов и подсистемой экспорта результатов.

Ключевой особенностью прототипа является сохранение методологической границы
между многоагентным обучением и экономическим механизмом распределения:
обучаемая политика управляет только операционными режимами периферийных узлов,
тогда как распределение задач, платежи и оценки внешних эффектов формируются
отдельным слоем назначения. Это делает программную реализацию согласованной с
теоретической частью диссертационной работы и подготавливает основу для
экспериментальной валидации, представленной в следующей главе.
